\hypertarget{index_Introduction}{}\section{Introduction}\label{index_Introduction}
Nous avons reçu comme projet de réaliser un interpréteur de grafcet. Nous allons vous présenter ici les différentes parties de ce projet. Vous trouverez dans le menu de gauche la description de chaque blocks. Pour réaliser ce projet nous avons programmé un compilateur, un assembleur et une machine virtuelle. Voici le fonctionnement de l\textquotesingle{}interpréteur \+: pour commencer le compilateur va analyser un grafcet et générer un fichier avec des instructions en assembleur puis l\textquotesingle{}assembleur va convertir le fichier assembleur en opcode. Enfin la machine virtuelle va être capable d\textquotesingle{}interpréter cet opcode et de générer de nouvelles valeurs de variable. Par la suite le projet a évolué il fallait être capable d\textquotesingle{}interpréter le grafcet sur le microcontroleur. Pour cela il falait envoyer l\textquotesingle{}opcode généré par l\textquotesingle{}assembleur sur le microcontôleur c\textquotesingle{}est le role du programme serialtransfer. La machine virtuelle se trouvera alors sur le microcontroleur et le tableau de variables sera constutué des entrées d\textquotesingle{}un système, comme par exemple des capteurs, et de ses sorties, par exemple un moteur.